% Generated draft: requires mathematical and source review.
\documentclass[11pt]{article}
\usepackage[margin=25mm]{geometry}
\usepackage{fontspec}
\setmainfont{Latin Modern Roman}
\usepackage{amsmath,amssymb,mathtools,unicode-math}
\setmathfont{Latin Modern Math}
\usepackage{enumitem,needspace,xurl,hyperref}
\setlength{\parindent}{0pt}
\setlength{\parskip}{5pt}
\setlength{\emergencystretch}{3em}
\newcommand{\N}{\mathbb{N}}
\newcommand{\Z}{\mathbb{Z}}
\newcommand{\Q}{\mathbb{Q}}
\newcommand{\R}{\mathbb{R}}
\newcommand{\C}{\mathbb{C}}
\newcommand{\F}{\mathbb{F}}
\newcommand{\A}{\mathbb{A}}
\newcommand{\PP}{\mathbb{P}}
\newcommand{\E}{\mathbb{E}}
\newcommand{\eps}{\varepsilon}
\newcommand{\dd}{\,\mathrm{d}}
\begin{document}
\section*{UnsolvedMath ID: 30005808}
\textbf{Source problem code:} catalogue entry
\subsection{Definitions and mathematical statement}
For a positive integer \(n\), write \([n]=\{1,\ldots,n\}\) and regard \([n]^2\) as a subset of the Euclidean plane \(\R^2\). For \(p\in[0,1]\), let \(R_{n,p}\subseteq[n]^2\) be random, with each point retained independently with probability \(p\). A set \(S\subseteq[n]^2\) is \emph{collinear-triple-free} if no three distinct points of \(S\) lie on a common affine line in \(\R^2\). Equivalently, for all pairwise distinct \((x_1,y_1),(x_2,y_2),(x_3,y_3)\in S\),
\[
\det\begin{pmatrix}
x_1&y_1&1\\
x_2&y_2&1\\
x_3&y_3&1
\end{pmatrix}\ne0.
\]
Define the random variable
\[
\alpha([n]^2,p)=\max\bigl\{|S|:S\subseteq R_{n,p},\ S\text{ is collinear-triple-free}\bigr\}.
\]
The empty set is admissible, so this maximum is defined for every realization, including when \(R_{n,p}=\varnothing\). The definition includes the boundary cases \(p=0\) and \(p=1\).

\textbf{Problem (source Problem 2).} Characterize \(\alpha([n]^2,p)\) as a function of the grid size \(n\) and retention probability \(p\).

\textbf{Unresolved scope.} The supplied source does not specify whether a characterization means an exact distribution, expectation, concentration, asymptotic order, or sharp asymptotic bounds. It also specifies neither a regime for \(p\), possibly depending on \(n\), nor a probability or error criterion for an asymptotic assertion. Thus the random variable is precisely defined here, but the requested characterization remains an open-ended research target.

The supplied report context motivates the grid problem by a finite-field analogue: it describes using a hypergraph container lemma and balanced supersaturation of collinear triples, obtained through pseudorandomness of the point--line incidence bipartite graph of the finite-field plane. These are contextual descriptions of the finite-field work, not hypotheses or asserted conclusions for the grid problem.
\subsection{Short English statement}
Independently keep each point of an \(n\)-by-\(n\) integer grid with probability \(p\). Characterize the largest number of retained points that can be selected with no three on a line. This maximum is random; the source leaves the desired type of characterization and parameter regime unspecified.
\subsection{Sources}
\begin{itemize}
\item Local registry: \texttt{data/unsolvedmath/problems.json}, permanent UnsolvedMath ID \(30005808\), title \emph{Collinear-Triple-Free Sets in Random Grids}, exact source code \texttt{catalogue entry}. The supplied record attributes the problem to the Oberwolfach Reports collection, gives proposed year 2024, and names no individual proposer.
\item \emph{Discrete Geometry} (2024), Oberwolfach Reports, pp.~137--202, DOI \texttt{10.4171/owr/2024/3}, source Problem 2. Supplied source and literature link: \url{https://doi.org/10.4171/owr/2024/3}.
\item Supplied literature assessment dated \texttt{2026-08-22}: the grid characterization is identified as a future direction, the finite-field analogue is described as treated, and no complete grid characterization is cited. The supplied status is \emph{open}. This is supplied dated curation, not an independently verified literature assessment; the record itself requests independent verification of time-sensitive status and cited claims.
\item The registry labels the statement \texttt{corrected\_verified} and reports that it was reconstructed and checked against the official Oberwolfach report. These are supplied verification claims, not independent verification by this draft.
\end{itemize}
\end{document}
